% ============================================================
%  Motivación
% ============================================================
\section{Motivación y planteo del problema}
\label{sec:motivacion}

\subsection{La necesidad de un sistema automatizado}

Los ensayos de seguridad eléctrica son esenciales para certificar y comercializar
productos eléctricos. Los métodos manuales presentan limitaciones como la
variabilidad de los resultados, mayores tiempos de prueba, riesgos operativos y
errores en la toma de datos. El desarrollo de un sistema automatizado de principio
a fin para realizar estos ensayos tiene como objetivo mejorar la precisión,
reducir la intervención humana y optimizar el tiempo de evaluación.

\subsection{Beneficios esperados}
\begin{itemize}
    \item Mayor precisión y repetibilidad en los resultados.
    \item Reducción de tiempos mediante la automatización de la secuencia de
          ensayos.
    \item Historial digital de los ensayos, que asegura la trazabilidad,
          facilita el análisis, respalda la certificación de productos y sirve de
          documentación en auditorías.
    \item Mayor cumplimiento normativo, al aplicar de forma sistemática los
          criterios de la norma.
\end{itemize}
